\subsection{Two-Dimensional Kolmogorov Flow}
\label{sec:results-kflow}

Our second testbed is two-dimensional (2D) Kolmogorov flow: the vorticity formulation of the incompressible Navier--Stokes equations on the periodic domain $[0,2\pi]^2$, with sinusoidal forcing at wavenumber $k=4$, viscosity $\nu=10^{-2}$, and linear (Rayleigh) drag $\alpha=0.1$. The state is a single vorticity field $\omega$ on a $64\times64$ grid ($4096$ degrees of freedom), advanced with a pseudo-spectral fourth-order Runge--Kutta (RK4) integrator at outer step $0.18$. The observation operator $\Hop$ subsamples every sixteenth point in each direction, so only a $4\times4=16$-point subset of the $4096$ grid points is observed. Unlike Lorenz-96, there is \emph{no} correlation preconditioning here: forming and factorizing a full $4096\times4096$ spatial covariance is prohibitively expensive, so we optimize the vorticity $\omega_0$ directly with an identity background, $\Bcov=\sigb^2\bm{I}$.

We emphasize at the outset that these are deliberately \emph{smoke-scale} validation runs that demonstrate the pipeline mechanics, not converged scientific results. The pseudo-spectral solver relies on \texttt{torch.fft.rfft2}, which is CPU-only in our stack, so every Kolmogorov experiment runs on CPU. The sliding-window ensemble is only $N=8$; the inverse observation operator is trained with lightweight defaults ($200$ trajectories, $30$ epochs) rather than the paper-scale $32{,}000$ trajectories and $500$ epochs; and the sliding-window experiments sweep only a single noise level, $\sigobs=0.0$. All quantitative claims below should be read with these caveats foregrounded.

\paragraph{Turbulent trajectory and warmup diagnostic.}
Figure~\ref{fig:kflow-traj} shows vorticity snapshots from a representative training trajectory. The fields display the characteristic turbulent shear bands at the $k=4$ forcing scale, with vorticity spanning roughly $[-12.25,\,14.08]$; this is the turbulent state that data assimilation (DA) must reconstruct from only a $4\times4$ subsampled observation grid. Warmed-up initial conditions are produced by integrating a spectrally-filtered random vorticity field forward $N_{\text{warmup}}=100$ outer steps into the statistically stationary regime. Figure~\ref{fig:kflow-warmup} tracks the enstrophy $\tfrac12\langle\omega^2\rangle$ along a single warmup trajectory: it starts at $0.500$ from the smooth cold start, overshoots to a peak near $15.3$ around step $20$ as small scales spin up, then relaxes into a quasi-stationary band, reading $9.824$ at $t=100$ and $6.195$ at $t=150$. The last-$20$-step relative drift is $10.82\%$, marginally above the $\lesssim10\%$ heuristic stationarity target, so we report the flow honestly as \emph{approximately} quasi-stationary rather than fully converged --- a direct consequence of the short single-trajectory diagnostic. The accompanying snapshots evolve from a featureless filtered field into developed $k=4$ turbulence by step $25$ onward.

\begin{figure}[!ht]
  \centering
  \includegraphics[width=0.82\linewidth]{KFlow_Full4DVar__cell06__out00.png}
  \caption{Vorticity snapshots of a single representative training trajectory for the 2D Kolmogorov flow on the $64\times64$ grid, shown at five equally spaced times along the $T=10$-step assimilation window. Each panel plots the vorticity $\omega$ with a diverging red--blue colormap (\texttt{RdBu\_r}) scaled symmetrically to $\pm\max|\omega|$. The fields exhibit the characteristic turbulent shear bands at the $k=4$ forcing scale that the assimilation experiments must reconstruct; the vorticity range across this trajectory is approximately $[-12.25,\,14.08]$. This is the turbulent, statistically stationary state being assimilated from only a $4\times4$ subsampled observation grid.}
  \label{fig:kflow-traj}
\end{figure}

\begin{figure}[!ht]
  \centering
  \includegraphics[width=0.82\linewidth]{KFlow_SlidingWindow__cell19__out00.png}
  \caption{Warmup and stationarity diagnostic for the Kolmogorov-flow spin-up, on a single trajectory. Left: enstrophy $\tfrac12\langle\omega^2\rangle$ versus outer integration step; the red dashed line marks the warmup length $N_{\text{warmup}}=100$. Enstrophy starts at $0.500$ from the smooth filtered cold start, overshoots to a peak near $15.3$ around step $20$, then relaxes to $9.824$ at $t=100$ and $6.195$ at $t=150$. The remaining panels show vorticity snapshots (\texttt{RdBu\_r}, shared scale) at $t=0,25,50,100,150$, evolving from a featureless field into developed $k=4$ turbulent shear bands. The last-$20$-step relative drift is $10.82\%$, marginally above the $\lesssim10\%$ heuristic, so the flow is reported as approximately quasi-stationary. Integrating $100$ warmup steps brings the cold start into a turbulent, approximately stationary regime suitable for generating assimilation experiments.}
  \label{fig:kflow-warmup}
\end{figure}

\paragraph{Experiment A: single $10$-step window.}
With $T_{\text{obs}}=\texttt{WINDOW\_T}=10$, cold-start background variance $\sigb=1.0$, and $\sigobs=0.0$, we run the four initialization$\times$optimization combinations (learned-inverse ``invobs'' or bicubic-interpolation ``baseline'' initialization, paired with hybrid or observation-only optimization). The clearest discriminator is the final limited-memory Broyden--Fletcher--Goldfarb--Shanno (L-BFGS) cost: invobs initialization lands the optimizer in a far lower-loss basin than the baseline. The final (unnormalized, summed $\Jb+\Jo$) losses are invobs+hybrid $56{,}385$ and invobs+obs-only $45{,}760$ versus baseline+hybrid $168{,}341$ and baseline+obs-only $138{,}254$ --- invobs starts the optimizer roughly $3\times$ lower in cost, at comparable iteration counts ($\approx261$--$270$). The analysis root-mean-square error (RMSE) at the window start orders the same way: the invobs combinations reach $\approx3.5$--$3.7$ while the baseline combinations sit near $4.8$--$5.5$ (Figure~\ref{fig:kflow-expA}). The spatial-RMSE-over-window curves separate clearly during assimilation, but all rollouts converge toward the saturation RMSE once the forecast leaves the window, as expected for a chaotic flow observed at only $16$ of $4096$ points.

\begin{figure}[!ht]
  \centering
  \includegraphics[width=0.82\linewidth]{KFlow_SlidingWindow__cell24__out00.png}
  \caption{Kolmogorov Experiment A: a single $10$-step assimilation window ($T_{\text{obs}}=\texttt{WINDOW\_T}=10$, $\sigb=1.0$, $\sigobs=0.0$), comparing the four initialization$\times$optimization combinations. Top: ensemble-mean spatial RMSE ($\pm1$ standard deviation over the $N=8$ ensemble) versus outer step, spanning the assimilation window (shaded, $0$--$10$, with the dashed window-end line) and the free forecast to step $25$. Bottom: bar chart of analysis RMSE at the window start ($t=0$). The invobs combinations reach lower in-window RMSE ($\approx3.5$--$3.7$) than the baseline combinations ($\approx4.8$--$5.5$), but all curves converge toward saturation once the forecast leaves the window, reflecting the chaotic flow and extreme observation sparsity ($16$ of $4096$ points). The final L-BFGS costs (invobs+hybrid $56{,}385$, invobs+obs-only $45{,}760$ versus baseline+hybrid $168{,}341$, baseline+obs-only $138{,}254$) show invobs initialization places the optimizer in a basin roughly $3\times$ lower in cost.}
  \label{fig:kflow-expA}
\end{figure}

\paragraph{Experiment B: stride sweep.}
We next run sliding-window (cycling) 4D-Var with invobs initialization every cycle, hybrid optimization, the propagated previous analysis as the $\Jb$ background on cycles $\geq1$, and a tighter cycling background variance $\sigb=0.3$. Figure~\ref{fig:kflow-expB} sweeps strides $\{2,5,10\}$ for total observation horizons $T_{\text{obs}}\in\{20,30\}$ at $\sigobs=0.0$. The stride curves overlap heavily within the shaded last-assimilation window and rise to the same saturation RMSE ($\approx4.6$--$4.8$) once the forecast begins; the coarsest stride, $10$, gives the lowest mean forecast RMSE for \emph{both} horizons, and the best single-window reference method is invobs+hybrid. The dashed best-Experiment-A reference, shifted to the last-window position, tracks the cycling curves closely, indicating that the cycling machinery does not yet separate from a single well-initialized window at this smoke scale.

\begin{figure}[!ht]
  \centering
  \includegraphics[width=0.82\linewidth]{KFlow_SlidingWindow__cell27__out00.png}
  \caption{Kolmogorov Experiment B1: sliding-window (cycling) 4D-Var stride sweep at $\sigobs=0.0$, with invobs initialization every cycle, hybrid optimization, the propagated previous analysis as the $\Jb$ background on cycles $\geq1$, and cycling $\sigb=0.3$. Rows are observation horizons $T_{\text{obs}}=20$ (top) and $30$ (bottom); only the $\sigobs=0.0$ column is populated (the right column of the $2\times2$ layout is empty because only one noise level was swept). Each panel plots ensemble-mean spatial RMSE ($\pm1$ standard deviation) versus outer step for strides $2$, $5$, and $10$ (coloured), over the last assimilation window (shaded, ending at the dashed line at $T_{\text{obs}}$) and the forecast to $T_{\text{obs}}+15$; the black dashed curve is the best single-window reference (invobs+hybrid) shifted to the last-window position. The stride curves overlap heavily and rise to the same saturation RMSE ($\approx4.6$--$4.8$); the coarsest stride $10$ gives the lowest mean forecast RMSE. At this smoke scale ($N=8$, one noise level) cycling does not yet separate from the single-window reference.}
  \label{fig:kflow-expB}
\end{figure}

\paragraph{Experiment C: cycling versus single-window.}
Figure~\ref{fig:kflow-expC} is the central comparison: the best-stride sliding-window run (full per-cycle analysis chain) against single-window invobs+hybrid and single-window baseline+obs-only, the latter two applied only to the last $10$ observations. The sliding-window-best curve traces per-cycle analysis RMSE from $t=0$ as a dashed marked chain and hands off to a solid curve over the last cycle plus its forecast. All three methods converge to the same saturation band ($\approx4.6$) once the forecast leaves the window. The honest reading of this smoke-scale run is a negative one: with $N=8$ and a single noise level, cycling does not deliver a decisive forecast advantage over a single well-initialized invobs+hybrid window once predictability is lost. (This contrasts with the Lorenz-96 cycling results of Section~\ref{sec:results-l96}, where a noise- and history-dependent cycling benefit was visible at $\sigobs=0.5$; the Kolmogorov runs sweep only $\sigobs=0$, the regime where cycling helps least, so the comparison here is not apples-to-apples.)

\begin{figure}[!ht]
  \centering
  \includegraphics[width=0.82\linewidth]{KFlow_SlidingWindow__cell31__out00.png}
  \caption{Kolmogorov Experiment C: the best-stride sliding-window run versus single-window baselines at $\sigobs=0.0$. Rows are $T_{\text{obs}}=20$ (top) and $30$ (bottom); only the $\sigobs=0.0$ column is populated. Each panel plots ensemble-mean spatial RMSE ($\pm1$ standard deviation) versus outer step. The sliding-window-best result is a dashed marked curve of per-cycle analysis RMSE from $t=0$, continued as a solid curve over the last cycle plus its forecast; it is compared against single-window invobs+hybrid and single-window baseline+obs-only, each applied to the last $10$ observations (shaded last-window region, dashed window-end line at $T_{\text{obs}}$, forecast to $T_{\text{obs}}+15$). All three methods converge to the same saturation band ($\approx4.6$) once the forecast leaves the window: at this smoke scale cycling delivers no decisive advantage over a single well-initialized invobs+hybrid window.}
  \label{fig:kflow-expC}
\end{figure}

\paragraph{Vorticity snapshot comparison.}
Finally, Figure~\ref{fig:kflow-snap} gives a qualitative $5\times3$ view of the reconstructed fields for Experiment C at $T_{\text{obs}}=20$ (ensemble sample $0$). The rows are ground-truth vorticity, the sliding-window-best prediction and its signed error, and the single-window invobs+hybrid prediction and its signed error; the columns are window start, window end, and forecast end. At window start, both methods reproduce the large-scale vorticity structures with small error. By forecast end, the error fields grow to the magnitude of the state itself, visually corroborating the RMSE saturation seen in the quantitative panels and underscoring the fundamental difficulty of constraining a turbulent $64\times64$ field from only $16$ point observations.

\begin{figure}[!ht]
  \centering
  \includegraphics[height=0.56\textheight,keepaspectratio]{KFlow_SlidingWindow__cell33__out00.png}
  \caption{Kolmogorov vorticity-field snapshot comparison for Experiment C at $T_{\text{obs}}=20$, $\sigobs=0.0$ (ensemble sample $0$). The $5\times3$ grid has columns at three times relative to the last-window start: window start ($t=10$), window end ($t=20$), and forecast end ($t=35$). Rows top to bottom: ground-truth vorticity; the sliding-window-best prediction; the sliding-window-best signed error; the single-window invobs+hybrid prediction; and the invobs+hybrid signed error. State rows use a diverging red--blue colormap (\texttt{RdBu\_r}) scaled to $\pm\max|\text{truth}|$; the two error rows share a single symmetric blue--white--red scale (\texttt{bwr}) showing prediction minus truth. At window start both methods reproduce the large-scale structures with small error, but by forecast end the error fields grow to the magnitude of the state itself, visually corroborating the RMSE saturation and the difficulty of reconstructing a $64\times64$ turbulent field from $16$ point observations.}
  \label{fig:kflow-snap}
\end{figure}

In summary, the Kolmogorov-flow runs reproduce the paper's central mechanism --- learned-inverse initialization places 4D-Var in a markedly lower-loss basin (Experiment~A) and produces the best in-window analyses --- but at this CPU-limited, $N=8$, single-noise-level smoke scale the operational cycling extension does not yet demonstrate a clear forecast advantage, and all methods saturate to climatological error once the chaotic flow decorrelates from truth. These outcomes are best understood as a working demonstration of the pipeline on a hard 2D system rather than as converged comparisons.
